% year: תשפ"ב
% type: מבחן
% semester: א
% moed: מיוחד
% number: 1
% points: 17
% categories: continuity, func-limits
% source: exams/מבחנים/תשפב/מועד ג תשפב.pdf

\begin{question}
עבור פרמטרים כלשהם $a,b\in\R$ נגדיר פונקציה על ידי
\[ f(x)=\begin{cases} e^{\frac1x}, & x<0,\\ ax+b-6, & 0\le x\le 9,\\ \frac{x^2-8x-9}{3-\sqrt{x}}, & x>9 \end{cases} \]
לאילו ערכים של הפרמטרים $a,b$ הפונקציה רציפה ב-$\R$? נמקו.
\end{question}

\begin{hint}
בכל אחד מהתחומים הפתוחים $f$ היא הרכבה/מנה של פונקציות רציפות. נשאר לבדוק את נקודות התפר $x=0$ ו-$x=9$.
\end{hint}

\begin{hint}
ב-$x=9$: $x^2-8x-9=(x-9)(x+1)$ ו-$x-9=(\sqrt x-3)(\sqrt x+3)$.
\end{hint}

\begin{solution}
\textbf{רציפות בתוך התחומים.} ב-$(-\infty,0)$ הפונקציה $e^{1/x}$ רציפה (הרכבת רציפות). ב-$(0,9)$ $f$ פולינום. ב-$(9,\infty)$ $f$ מנה של פונקציות רציפות שהמכנה שלה $3-\sqrt x<0$ אינו מתאפס. לכן $f$ רציפה בכל $x\neq0,9$ לכל $a,b$, ויש לבדוק רק את $x=0$ ואת $x=9$.

\textbf{הנקודה $x=0$.} $f(0)=b-6$, והגבול מימין הוא $\lim_{x\to0^+}(ax+b-6)=b-6$. משמאל: כאשר $x\to0^-$ מתקיים $\frac1x\to-\infty$, ולכן
\[ \lim_{x\to0^-}e^{1/x}=0. \]
רציפות ב-$0$ שקולה ל-$b-6=0$, כלומר $b=6$.

\textbf{הנקודה $x=9$.} $f(9)=9a+b-6$, וזה גם הגבול משמאל. מימין, עבור $x>9$:
\[ \frac{x^2-8x-9}{3-\sqrt x}=\frac{(x-9)(x+1)}{3-\sqrt x}=\frac{(\sqrt x-3)(\sqrt x+3)(x+1)}{-(\sqrt x-3)}=-(\sqrt x+3)(x+1), \]
ולכן
\[ \lim_{x\to9^+}f(x)=-(3+3)(9+1)=-60. \]
רציפות ב-$9$ שקולה ל-$9a+b-6=-60$.

\textbf{פתרון המערכת.} מ-$b=6$ נקבל $9a=-60$, כלומר $a=-\frac{20}{3}$.

\textbf{תשובה:} $f$ רציפה בכל $\R$ אם ורק אם $\boxed{a=-\tfrac{20}{3},\ b=6}$.
\end{solution}
